\section{Teorias de primeira ordem}\label{sec:teoriasPrimeiraOrdem}
\index{teoria da prova}
\index{sistema de Hilbert}
Até aqui, discutimos o que é uma linguagem de primeira ordem e aprendemos a identificar termos, sentenças e fórmulas.
O próximo passo é entender teorias, teoremas e demonstrações de teorias de primeira ordem.

A apresentação exata dos axiomas lógicos adotados varia muito de texto para texto.
Há quem diga que há tantas formas de axiomatizar a lógica quanto há textos de lógica.
Acreditamos que, em algum sentido, este texto somará 1 a essa contagem, mesmo sem apresentar nenhuma inovação.
Não comentaremos outras formas de axiomatização com profundidade.

Agora estamos prontos para enunciar os axiomas lógicos adotados neste texto.
Como fizemos com a lógica proposicional, adotamos um sistema de axiomas estilo Hilbert.
Ressaltamos novamente que esse não é o único modo de formalizar a dedução lógica.

\begin{metadefinition}[Axiomas da lógica de primeira ordem]\label{mdef:axiomasLogicaPrimeiraOrdem}\index{axioma!lógica de primeira ordem}
    Fixado um léxico $\mathcal L$ para uma linguagem de primeira ordem, os axiomas da Lógica de Primeira Ordem em $\mathcal L$ são as fórmulas que podem ser escritas de alguma das seguintes formas, em que $\phi$, $\psi$ e $\chi$ são $\mathcal L$-fórmulas quaisquer, $x$, $x_1,\dots,x_n$, $y_1,\dots,y_n$ são variáveis, $t$ é um $\mathcal L$-termo, e $P$ e $f$ são símbolos de predicado e de função, respectivamente.
    \begin{enumerate}[label=(FO\arabic*), leftmargin=3.2em]
        \item $\phi \rightarrow (\psi \rightarrow \phi)$;
        \item $(\phi \rightarrow (\psi \rightarrow \chi)) \rightarrow ((\phi \rightarrow \psi) \rightarrow (\phi \rightarrow \chi))$;
        \item $(\phi \rightarrow \psi) \rightarrow ((\phi \rightarrow \neg \psi) \rightarrow \neg \phi)$;
        \item $(\phi \land \psi) \rightarrow \phi$;
        \item $(\phi \land \psi) \rightarrow \psi$;
        \item $\phi \rightarrow (\psi \rightarrow (\phi \land \psi))$;
        \item $\phi \rightarrow (\phi \lor \psi)$;
        \item $\psi \rightarrow (\phi \lor \psi)$;
        \item $(\phi \rightarrow \chi) \rightarrow ((\psi \rightarrow \chi) \rightarrow ((\phi \lor \psi) \rightarrow \chi))$;
        \item $(\psi \leftrightarrow \phi) \rightarrow (\psi \rightarrow \phi)$;
        \item $(\psi \leftrightarrow \phi) \rightarrow (\phi \rightarrow \psi)$;
        \item $(\psi \rightarrow \phi) \rightarrow ((\phi \rightarrow \psi) \leftrightarrow (\psi \leftrightarrow \phi))$;
        \item $\neg \neg \phi \rightarrow \phi$;
        \item $\forall x \phi \rightarrow \phi[x/t]$, se $x$ é substituível por $t$ em $\phi$;
        \item $\phi[x/t] \rightarrow \exists x \phi$, se $x$ é substituível por $t$ em $\phi$;
        \item $x=x$;
        \item $\bigl(\bigwedge_{i=1}^n x_i=y_i\bigr) \rightarrow (f(x_1, \dots, x_n)=f(y_1, \dots, y_n))$, para cada símbolo funcional $f$ de aridade $n>0$.
        \item $\bigl(\bigwedge_{i=1}^n x_i=y_i\bigr) \rightarrow (P(x_1, \dots, x_n) \leftrightarrow P(y_1, \dots, y_n))$, em que $P$ é um símbolo de predicado não lógico de aridade $n>0$, ou o símbolo lógico de igualdade $=$ (que possui aridade 2).
        \end{enumerate}

        Observação: quando $n=1$, $\bigwedge_{i=1}^n\phi_i$ é $\phi_1$.
        Recursivamente para $n>1$, $\bigwedge_{i=1}^n\phi_i$ é a fórmula $\Big(\bigwedge_{i=1}^{n-1}\phi_i\Big)\land\phi_n$.
\end{metadefinition}

Os axiomas (FO1)--(FO13) possuem a mesma forma dos axiomas proposicionais (PC1)--(PC13), agora permitindo que $\phi$, $\psi$ e $\chi$ sejam fórmulas de primeira ordem.

Os axiomas (FO14) e (FO15) governam o comportamento dos quantificadores em relação à substituição de termos por variáveis.

Por fim, (FO16)--(FO18) são os axiomas que governam a igualdade e sua compatibilidade com símbolos funcionais e relacionais.

A seguir, definiremos as noções de teoria, demonstração e teorema.

\begin{metadefinition}[Teoria]\label{mdef:teoriaPrimeiraOrdem}\index{teoria}\index{axioma!não lógico}
    Uma \emph{teoria de primeira ordem} consiste em um par $T = (\mathcal L, \Sigma)$, em que $\mathcal L$ é um léxico de primeira ordem e $\Sigma$ é uma coleção de $\mathcal L$-sentenças, chamadas de \emph{axiomas não lógicos} de $T$, ou simplesmente de \emph{axiomas} de $T$.
\end{metadefinition}

\begin{metadefinition}[Demonstração]\label{mdef:demonstracaoFormal}\index{demonstração formal}\index{teorema}\index{modus ponens}
    Seja $T=(\mathcal L,\Sigma)$ uma teoria de primeira ordem e seja $\phi$ uma $\mathcal L$-fórmula.
    Uma \emph{demonstração} de $\phi$ a partir de $T$ é uma sequência finita de $\mathcal L$-fórmulas $\psi_0, \dots, \psi_n$ tal que $\psi_n$ é $\phi$ e, para cada $i\leq n$:
    \begin{enumerate}[label=(\roman*)]
        \item $\psi_i$ é um axioma lógico, ou
        \item $\psi_i$ é um axioma de $T$, ou
        \item $\psi_i$ é obtida por \emph{modus ponens} a partir de duas fórmulas anteriores.
        Ou seja, existem $j, k<i$ tais que $\psi_j$ é  $\psi_k \rightarrow \psi_i$.
        \item Existe uma variável $x$, fórmulas $\theta$ e $\eta$, e $j<i$, tais que $\psi_j$ é $\theta\rightarrow \eta$ e $\psi_i$ é $\theta\rightarrow \forall x \eta$, em que $x$ não ocorre livre em $\theta$.
        Essa regra é chamada de introdução de quantificador universal.
        \item Existe uma variável $x$, fórmulas $\theta$ e $\eta$, e $j<i$, tais que $\psi_j$ é $\theta\rightarrow \eta$ e $\psi_i$ é $\exists x \theta\rightarrow \eta$, em que $x$ não ocorre livre em $\eta$.
        Essa regra é chamada de introdução de quantificador existencial.
    \end{enumerate}

    Caso exista uma demonstração de $\phi$ a partir de $T$, dizemos que $\phi$ é um \emph{teorema} de $T$, e escrevemos $T \vdash \phi$.
    Sendo $T=(\mathcal L, \Sigma)$, também podemos escrever $\Sigma \vdash_\mathcal L \phi$.

    Se tal demonstração só usa, além dos axiomas de $T$, os axiomas proposicionais de primeira ordem (FO1)--(FO13) e a regra de \emph{modus ponens}, dizemos que $\phi$ é uma \emph{consequência proposicional} de $T$, e escrevemos $T \vdash_P \phi$.
\end{metadefinition}

Na prática, raramente escrevemos, em matemática, uma demonstração tão formal quanto a Metadefinição~\ref{mdef:demonstracaoFormal} exige.
Para reforçar a ideia de que as demonstrações usuais em matemática poderiam, em tese, ser reescritas de forma tão formal, apresentaremos alguns resultados relevantes.

Começaremos ``importando'' teoremas da lógica proposicional para a lógica de primeira ordem.

A definição abaixo é similar àquela apresentada em Metadefinição~\ref{mdef:substituicaoProposicional} para a lógica proposicional, mas agora permite substituir letras proposicionais por fórmulas de uma linguagem de primeira ordem.

\begin{metadefinition}
    A \emph{extensão da lógica proposicional para a lógica de primeira ordem} é a teoria $T_P=(\mathcal L_P,\varnothing)$, em que os únicos símbolos não lógicos de $\mathcal L_P$ são letras proposicionais, isto é, símbolos de predicado $0$-ários $P_0$, $P_1$, $P_2,\dots$.

    Exceto quando explicitamente mencionado, vamos assumir que qualquer outra teoria que estudaremos não possui as letras proposicionais $P_0$, $P_1$, $P_2$, $\dots$ como símbolos de seu léxico.
\end{metadefinition}

Note que as fórmulas proposicionais são exatamente as $\mathcal L_P$-fórmulas sem quantificadores e sem símbolos de igualdade.
\begin{metadefinition}
   Seja $\mathcal L$ um léxico de primeira ordem.
   Considere uma coleção finita $\mathcal P$ de letras proposicionais de $\mathcal L_P$.
   Uma atribuição de $\mathcal P$ a $\mathcal L$-fórmulas é um mapeamento tal que, para cada letra proposicional $P$ em $\mathcal P$, $\rho(P)$ é uma $\mathcal L$-fórmula.

   Se $\theta$ é uma $\mathcal L_P$-fórmula, dizemos que $\rho$ é uma atribuição para $\theta$ se toda letra proposicional que ocorre em $\theta$ está em $\mathcal P$.
\end{metadefinition}

Essa substituição é literal: ao substituir uma letra proposicional por uma fórmula, variáveis livres da fórmula substituída podem passar a ficar no escopo de quantificadores já presentes em $\theta$.

\begin{metadefinition}[Substituição proposicional]\label{mdef:substituicaoProposicionalFO}\index{substituição!proposicional}
   Seja $\mathcal L$ um léxico de primeira ordem.
   Considere uma coleção finita $\mathcal P$ de letras proposicionais de $\mathcal L_P$ e $\rho$ uma atribuição de $\mathcal P$ a $\mathcal L$-fórmulas.

   Recursivamente na complexidade de uma $\mathcal L_P$-fórmula $\theta$, definimos a substituição de $\theta$ via $\rho$, denotada por $\subst(\theta, \rho)$, a partir do seguinte algoritmo.
   Se $\rho$ não é uma atribuição para $\theta$, então $\subst(\theta, \rho)$ é um símbolo $I$ que não faz parte da linguagem (indicando, intuitivamente, que $\subst(\theta, \rho)$ é indefinido).
   Caso contrário, procedemos como nos itens abaixo.
   Note que, como os únicos símbolos não lógicos de $\mathcal L_P$ são predicados $0$-ários, os únicos termos de $\mathcal L_P$ são variáveis.
    \begin{itemize}
        \item Se $\theta$ é uma letra proposicional $P$, então $\subst(\theta, \rho)$ é $\rho(P)$.
        \item Se $\theta$ é $\neg \psi$, então $\subst(\theta, \rho)$ é $\neg \subst(\psi, \rho)$.
        \item Se $\theta$ é $x=y$, em que $x$ e $y$ são variáveis, então $\subst(\theta, \rho)$ é $x=y$.
        \item Se $\theta$ é $\psi \square \eta$, em que $\square$ é $\vee$, $\wedge$, $\rightarrow$ ou $\leftrightarrow$, então $\subst(\theta, \rho)$ é $\subst(\psi, \rho) \square \subst(\eta, \rho)$.
        \item Se $\theta$ é $\forall x \psi$ ou $\exists x \psi$, então $\subst(\theta, \rho)$ é $\forall x \subst(\psi, \rho)$ ou $\exists x \subst(\psi, \rho)$, respectivamente.
    \end{itemize}

    Se $\theta$ é uma tautologia proposicional e $\rho$ é uma atribuição para $\theta$, então $\subst(\theta, \rho)$ é dita uma \emph{instância de tautologia}.
\end{metadefinition}

\begin{metatheorem}
    Seja $\mathcal L$ um léxico de primeira ordem e $\phi$ uma $\mathcal L$-fórmula.
    Se $\phi$ é uma instância de tautologia, então, para toda teoria de primeira ordem $T$ de léxico $\mathcal L$,
    \[
        T\vdash_P \phi.
    \]
\end{metatheorem}

\begin{proof}
    Seja $\theta$ uma tautologia proposicional, seja $\mathcal P$ uma coleção finita de letras proposicionais e seja $\rho$ uma atribuição de $\mathcal P$ a $\mathcal L$-fórmulas tal que $\subst(\theta, \rho)$ é $\phi$.

    Como $\theta$ é uma tautologia, $\theta$ é um teorema da lógica proposicional, ou seja, $\vdash_P \theta$.

    Assim, existe uma demonstração proposicional de $\theta$ a partir da coleção vazia de axiomas, digamos $\theta_0, \dots, \theta_n$, em que $\theta_n$ é $\theta$.

    Seja $\mathcal P'$ uma coleção finita de letras proposicionais que contém $\mathcal P$ e todas as letras proposicionais que ocorrem em alguma das fórmulas $\theta_0,\dots,\theta_n$.

    Estenda $\rho$ a uma atribuição $\rho'$ de $\mathcal P'$ a $\mathcal L$-fórmulas, atribuindo arbitrariamente $\mathcal L$-fórmulas às letras de $\mathcal P'\setminus\mathcal P$.
    Assim, $\rho'$ é uma atribuição para cada $\theta_i$ e coincide com $\rho$ em todas as letras proposicionais que ocorrem em $\theta$.

    Para cada $i\leq n$, seja $\psi_i$ a fórmula $\subst(\theta_i, \rho')$.
    Note que $\psi_n$ é $\subst(\theta_n, \rho')$, que é $\subst(\theta, \rho)$, que é $\phi$.

    Veremos que $\psi_0, \dots, \psi_n$ é uma demonstração de $\phi$ que usa apenas os axiomas proposicionais de primeira ordem (FO1)--(FO13) e a regra de \emph{modus ponens}.

    Com efeito, para cada $i\leq n$, $\theta_i$ é um axioma da lógica proposicional ou é obtida por \emph{modus ponens} de duas fórmulas anteriores.
    Se $\theta_i$ é um axioma proposicional do tipo (PC1)--(PC13), então $\psi_i$ é um axioma lógico do tipo (FO1)--(FO13) (mesma numeração).

    Se $\theta_i$ é obtida por \emph{modus ponens} de $\theta_j$ e $\theta_k$, com $j, k<i$ e $\theta_j$ é $\theta_k \rightarrow \theta_i$, então $\psi_j$ é $\psi_k \rightarrow \psi_i$, e, portanto, $\psi_i$ é obtida por \emph{modus ponens} de $\psi_j$ e $\psi_k$.

    Assim, para toda teoria $T$ de léxico $\mathcal L$, temos $T\vdash_P \phi$.
\end{proof}

O resultado a seguir pode ser interpretado como a asserção de que resultados já provados podem ser usados para provar novos resultados.
\begin{metalemma}[Modus ponens]\label{mlem:modusPonens}
    Seja $T$ uma teoria de primeira ordem e sejam $\phi$ e $\psi$ $\mathcal L$-fórmulas, em que $\mathcal L$ é o léxico de $T$.
    Se $T \vdash \phi$ e $T \vdash \phi \rightarrow \psi$, então $T \vdash \psi$.
\end{metalemma}
\begin{proof}
    Existem demonstrações de $\phi$ e de $\phi \rightarrow \psi$ a partir de $T$, digamos $\psi_0, \dots, \psi_n$ e $\psi'_0, \dots, \psi'_m$, respectivamente, de modo que $\psi_n$ é $\phi$ e $\psi'_m$ é $\phi \rightarrow \psi$.
    Considere a seguinte sequência de fórmulas:
    \begin{equation*}
        \psi_0, \dots, \psi_n, \psi'_0, \dots, \psi'_{m}, \psi.
    \end{equation*}

    A sequência $\psi_0, \dots, \psi_n, \psi'_0, \dots, \psi'_{m}$, por ser uma concatenação de duas demonstrações, satisfaz claramente as condições para ser uma demonstração.
    Ao acrescentar $\psi$ ao final da sequência, ainda temos as condições satisfeitas, uma vez que $\psi$ é obtida por \emph{modus ponens} a partir de $\phi$ e $\phi \rightarrow \psi$, que são, respectivamente, $\psi_n$ e $\psi'_m$.
\end{proof}

O resultado a seguir nos diz que se uma teoria $T'$ prova todos os axiomas de $T$, então $T'$ prova todos os teoremas de $T$.

\begin{metaproposition}\label{mprop:extensaoDeTeoria}
    Sejam $T=(\mathcal L,\Sigma)$ e $T'=(\mathcal L',\Sigma')$ teorias de primeira ordem tais que todo símbolo de $\mathcal L$ é um símbolo de $\mathcal L'$.
    Suponha que, para cada $\sigma\in\Sigma$, temos $T'\vdash \sigma$.

    Então, para cada $\mathcal L$-fórmula $\phi$, se $T\vdash \phi$, então $T'\vdash \phi$.

    Em particular, se $T'$ é obtida de $T$ acrescentando novos axiomas, então todo teorema de $T$ é um teorema de $T'$.
\end{metaproposition}

\begin{proof}
    Seja $\phi$ uma $\mathcal L$-fórmula tal que $T\vdash \phi$.
    Fixe uma demonstração $\phi_0,\dots,\phi_n$ de $\phi$ a partir de $T$.

    Provaremos, por indução em $i\leq n$, que $T'\vdash \phi_i$.

    Suponha que $T'\vdash \phi_j$ para todo $j<i$.
    Se $\phi_i$ é um axioma lógico, então $\phi_i$ também é um axioma lógico na linguagem $\mathcal L'$.
    Logo, $T'\vdash \phi_i$.

    Se $\phi_i$ é um axioma de $T$, então $\phi_i\in\Sigma$.
    Pela hipótese, $T'\vdash \phi_i$.

    Se $\phi_i$ é obtida por \emph{modus ponens} a partir de fórmulas anteriores, então a conclusão segue da hipótese indutiva e do Metalema~\ref{mlem:modusPonens}.

    Se $\phi_i$ é obtida por uma das regras de introdução de quantificadores a partir de uma fórmula anterior, então, pela hipótese indutiva, essa fórmula anterior é demonstrável em $T'$.
    Aplicando a mesma regra de introdução de quantificador, obtemos $T'\vdash \phi_i$.

    Portanto, $T'\vdash \phi_i$ para todo $i\leq n$.
    Como $\phi_n$ é $\phi$, segue que $T'\vdash \phi$.
\end{proof}

Os resultados acima afirmam que ``existem demonstrações''.
É importante observar que tais provas existenciais são construtivas: elas de fato fornecem um algoritmo para, dados os objetos declarados nas hipóteses, efetivamente redigir uma demonstração.

Abaixo, provaremos, como exemplo, um teorema da lógica pura que intuitivamente diz que o universo de discurso é não vazio.

\begin{metaproposition}\label{mprop:universoNaoVazio}
    Para cada teoria de primeira ordem $T$, $T \vdash \exists x (x=x)$.
\end{metaproposition}
\begin{proof}
    Considere a seguinte demonstração.
    \begin{align*}
    1.\ & x=x
        & \text{(FO16)} \\
    2.\ & x=x \rightarrow \exists x(x=x)
        & \text{(FO15, aplicado à fórmula $x=x$ e ao termo $x$)} \\
    3.\ & \exists x(x=x)
        & \text{(modus ponens de 1. e 2.)}
    \end{align*}
\end{proof}

\subsection*{Exercícios}

\begin{exercise}
    Explique por que toda instância de (FO1)--(FO13) é uma instância de tautologia.
\end{exercise}
